\subsection{Defense-in-Depth：原则、机制与模型 Hardening}

slides 64--75 从风险转入防御。总体原则是 defense-in-depth、least privilege、privilege separation、safe-by-design 和 secure-by-design。机制包括模型 hardening、输入 sanitization、动作 policy enforcement、身份/访问控制、任务分解与隔离、runtime monitoring、information-flow control 和 provenance。模型 alignment 是一层，但不能承担全部系统保证。

\lecturefigure{slides-images/slide-064.png}{官方 slide 64：Agentic Hybrid Systems 与安全挑战总结}
\lecturefigure{slides-images/slide-065.png}{官方 slide 65：Defense-in-depth 与 least privilege}
\lecturefigure{slides-images/slide-067.png}{官方 slide 67：Safe-by-design 与 secure-by-design 的完整原则}
\lecturefigure{slides-images/slide-068.png}{官方 slide 68：Defense mechanisms 总览起点}
\lecturefigure{slides-images/slide-070.png}{官方 slide 70：Defense mechanisms 的完整框架}
\lecturefigure{slides-images/slide-071.png}{官方 slide 71：AI Model Hardening 与 Alignment}
\lecturefigure{slides-images/slide-072.png}{官方 slide 72：从模型层继续向系统层布防}
\lecturefigure{slides-images/slide-075.png}{官方 slide 75：Least privilege 与 policy enforcement 的位置}

\begin{knowledgebox}{术语消化：五个设计原则}
\begin{tabular}{p{3.0cm}p{4.2cm}p{5.0cm}}
\toprule
\textbf{原则} & \textbf{核心问题} & \textbf{Agent 实现} \\
\midrule
Defense-in-depth & 单层失守怎么办 & 模型、策略、工具、业务系统独立校验 \\
Least privilege & 最少需要什么权限 & 按任务、时间和资源发放 capability \\
Privilege separation & 高风险能力如何隔离 & 拆进独立进程/Agent，显式审批通信 \\
Safe-by-design & 错误如何限制后果 & 可逆动作、幂等、默认拒绝、人工闸门 \\
Secure-by-design & 攻击者如何被约束 & threat model、最小接口、审计与更新 \\
\bottomrule
\end{tabular}
\end{knowledgebox}

\begin{warningbox}{模型更稳健不等于系统有边界}
即使 jailbreak rate 降低，Agent 仍可能因正常语言误解而执行错误动作。确定性 policy、typed tool、resource ACL 和 transaction invariant 必须独立于模型；这样模型层退化或更换时，系统约束仍存在。
\end{warningbox}

